\section{Sparse / Merge}
\subsection{Sparse Table}
Max/Min/GCD Query
\begin{minted}{c++}
struct SparseTable {
  SparseTable(vector<int>& data)
  {
    int n = data.size();
    int maxLog = log2(n) + 1;
    table.assign(n, vector<int>(maxLog));
    for (int i = 0; i < n; ++i) {
      table[i][0] = data[i];
    }
    for (int j = 1; (1 << j) <= n; ++j) {
      for (int i = 0; (i + (1 << j) - 1) < n; ++i) {
        table[i][j] = max(table[i][j - 1], table[i +
        (1 << (j - 1))][j - 1]);
      }
    }
  }
  int query(int L, int R) {
    int j = log2(R - L + 1);
    return max(table[L][j], table[R - (1 << j) +
    1][j]);
  }
  vector<vector<int>> table;
};
\end{minted}
\subsection{Sparse Table 2D}
\begin{minted}{c++}
struct SparseTable2D {
  static const int LG = 10;
  int n, m;
  vector<vector<int>> a;
  vector<int> lg2;
  vector<vector<vector<vector<int>>>> st;
  SparseTable2D(int n, int m) : n(n), m(m) {
    a.assign(n, vector<int>(m, 0));
    lg2.resize(max(n, m) + 1, 0);
    for (int i = 2; i <= max(n, m); i++) lg2[i] =
    lg2[i >> 1] + 1;
    st.assign(n, vector<vector<vector<int>>>(m,
    vector<vector<int>>(LG, vector<int>(LG,
    0))));
  }
  void set_val(int i, int j, int val) { a[i][j] = val; }
  void build() {
    for (int i = 0; i < n; i++)
      for (int j = 0; j < m; j++)
        st[i][j][0][0] = a[i][j];
    for (int a = 0; a < LG; a++) {
      for (int b = 0; b < LG; b++) {
        if (a + b == 0) continue;
        for (int i = 0; i + (1 << a) <= n; i++) {
          for (int j = 0; j + (1 << b) <= m; j++) {
            if (!a)
              st[i][j][a][b] = max(st[i][j][a][b-1],
              st[i][j + (1 << (b - 1))][a][b - 1]);
            else
              st[i][j][a][b] = max(st[i][j][a-1][b],
              st[i + (1 << (a - 1))][j][a - 1][b]);
          }
        }
      }
    }
  }
  int query(int x1, int y1, int x2, int y2) {
    int a = lg2[x2 - x1 + 1], b = lg2[y2 - y1 + 1];
    int m1 = max(st[x1][y1][a][b], st[x2 - (1 << a)
    + 1][y1][a][b]);
    int m2 = max(st[x1][y2 - (1 << b) + 1][a][b],
    st[x2 - (1 << a) + 1][y2 - (1 << b) + 1][a][b]);
    return max(m1, m2);
  }
};
\end{minted}
\begin{minted}{c++}
// st.set_val(i, j, grid[i][j]);
// st.build(); // st.query(0, 0, 1, 2); // max in
// submatrix [(0,0) to (1,2)]
\end{minted}
\subsection{Merge Sort Tree (Values \texorpdfstring{$\le$}{<=} K in Range)}
\begin{minted}{c++}
struct MergeSortTree {
  int n;
  vector<vector<int>> tree;
  MergeSortTree(const vector<int>& arr) {
    n = arr.size();
    tree.resize(4 * n);
    build(1, 0, n - 1, arr);
  }
  void build(int node, int l, int r, const vector<int>& arr) {
    if (l == r) {
      tree[node].push_back(arr[l]);
      return;
    }
    int mid = (l + r) / 2;
    build(2 * node, l, mid, arr);
    build(2 * node + 1, mid + 1, r, arr);
    merge(tree[2 * node].begin(), tree[2 * node].end(),
          tree[2 * node + 1].begin(), tree[2 * node + 1].end(),
          back_inserter(tree[node]));
  }
  // Count elements <= k in range [ql, qr]
  int query(int node, int l, int r, int ql, int qr, int k) {
    if (l > qr || r < ql) return 0;
    if (l >= ql && r <= qr) {
      return upper_bound(tree[node].begin(), tree[node].end(), k) - tree[node].begin();
    }
    int mid = (l + r) / 2;
    return query(2 * node, l, mid, ql, qr, k) +
           query(2 * node + 1, mid + 1, r, ql, qr, k);
  }
};
// To query distinct elements in [L, R]:
// 1. Compute next_right[i] for all elements.
// 2. Build MergeSortTree on next_right array.
// 3. Ans = (R - L + 1) - query(1, 0, n - 1, L, R, R)
\end{minted}
\subsection{K-th Smallest in Range (Index SegTree)}
\begin{minted}{c++}
struct KthSmallest {
  int n;
  vector<vector<int>> seg;
  KthSmallest(const vector<int>& arr) {
    n = arr.size();
    seg.resize(4 * n);
    vector<pair<int, int>> B(n);
    for (int i = 0; i < n; i++) B[i] = {arr[i], i};
    sort(B.begin(), B.end());
    build(0, 0, n - 1, B);
  }
  void build(int node, int l, int r, const vector<pair<int, int>>& B) {
    if (l == r) {
      seg[node].push_back(B[l].second);
      return;
    }
    int mid = (l + r) / 2;
    build(2 * node + 1, l, mid, B);
    build(2 * node + 2, mid + 1, r, B);
    merge(seg[2 * node + 1].begin(), seg[2 * node + 1].end(),
          seg[2 * node + 2].begin(), seg[2 * node + 2].end(),
          back_inserter(seg[node]));
  }
  // Returns the ORIGINAL array index of the K-th smallest element
  int query(int node, int st, int end, int l, int r, int k) {
    if (st == end) return seg[node][0];
    int p = upper_bound(seg[2 * node + 1].begin(), seg[2 * node + 1].end(), r)
          - lower_bound(seg[2 * node + 1].begin(), seg[2 * node + 1].end(), l);
    int mid = (st + end) / 2;
    if (p >= k) return query(2 * node + 1, st, mid, l, r, k);
    else return query(2 * node + 2, mid + 1, end, l, r, k - p);
  }
};
\end{minted}
